\section{Two Optimal Tiling Examples: V1 LP Solution at a Glance}
\label{app:tiling}
This section shows examples of optimal tiling, where accepting a locally slower/smaller tiling to enable better fusion or placement.
This showcases the strength of the CHEETAH-V1 formulation, which discovers nontrivial tiling-placement-fusion trade-offs that less expressive methods would miss. The first example is Llama-3 70B on a modeled H100; the second is DeepSeek-V3 on a modeled RTX~5090.

By fixing the architecture, the inner loop derives optimal tiling and scheduling strategies that can act as custom speedups for targeted inference tasks. Therefore this section briefly provides motivation for future work on expanded hardware variants. This is motivated and similar in spirit to the speedups seen be FlashAttention \citep{dao2022flashattention} due to tiling optimization. 

% \begin{table}[ht!]
% \centering
% \caption{Per-chip V1 LP scheduling lift over naive (greedy) scheduling
% on the same modelled hardware. Speedup $=$ \texttt{static.objective\_cycles}
% $\div$ \texttt{v1\_rounded.objective\_cycles}.
% Source: \texttt{exp3\_<hw>\_*.csv}.}
% \label{tab:v1_per_chip_lift}
% \small
% \begin{tabular}{@{}llr@{}}
% \toprule
% \textbf{Workload} & \textbf{Hardware} &
% \textbf{V1 LP speedup over naive} \\
% \midrule
% Llama-3 70B  & H100      & $3.06\times$ \\
% DeepSeek-V3  & RTX~5090  & $5.51\times$ \\
% \bottomrule
% \end{tabular}
% \end{table}

\subsection{V1 LP Solution: Llama-3 70B on modeled H100}

\noindent\textbf{Problem dimensions:}
$L = 720$ layers, $C = 20$ tile configs, $T = 720$ timesteps,
$E = 703$ activation edges.
Total LP variables: 1{,}601{,}280.
Objective: $542{,}790{,}817$ cycles.
\texttt{CONVERGED\_STRICT} in 1{,}800 PDLP iterations.

\paragraph{Tile-Config Selection ($y_{i,c}$):} The V1 LP heavily favours small/fast configurations. Table~\ref{tab:tile_band} groups tiles into three size bands, and Table~\ref{tab:config_histogram} gives the per-config histogram.

\begin{table}[h]
\centering
\caption{Tile-config band distribution across $L=720$ layers.}
\label{tab:tile_band}
\begin{tabular}{lrr}
\toprule
\textbf{Tile size band} & \textbf{Layers} & \textbf{\% of model} \\
\midrule
Small ($c = 0\ldots4$): fast / large working set  & 427 & 59\% \\
Medium ($c = 5\ldots14$)                            & 275 & 38\% \\
Large ($c = 15\ldots19$): slow / small working set &  18 &  2\% \\
\bottomrule
\end{tabular}
\end{table}

\noindent Config histogram (picks per config index):

\begin{table}[h]
\centering
\caption{Number of layers assigned to each tile config $c$.
Configs with zero picks are omitted.}
\label{tab:config_histogram}
\begin{tabular}{lrrrrrrrrrr}
\toprule
$c$ & 0 & 1 & 3 & 4 & 5 & 8 & 12 & 16 & 17 \\
\midrule
Layers & 285 & 1 & 1 & 140 & 1 & 148 & 125 & 15 & 3 \\
\bottomrule
\end{tabular}
\end{table}

\noindent Config $c=0$ (largest tile) was selected for ${\approx}40\%$
of layers, with $c \in \{4, 8, 12\}$ covering most of the remainder.
The LP decisively favours large tiles for this 70B model on H100:
with 1.05M MACs and a 50 MB global buffer, reducing tile size to fit
memory is unnecessary and incurs iteration overhead.

\paragraph{Per-Layer Execution Time $T_i$:}

\[
    T_i^{\min} = 1{,}204 \text{ cycles}, \quad
    T_i^{\mathrm{med}} = 524{,}288, \quad
    T_i^{\max} = 1{,}835{,}008, \quad
    \textstyle\sum T_i = 542{,}790{,}817.
\]

\noindent The ten longest layers (Table~\ref{tab:longest_layers}) all hit the same $1{,}835{,}008$-cycle bound: these are the large FFN gate/up/down and Q/K/V/output projection sublayers in the transformer middle blocks:

\begin{table}[h]
\centering
\caption{Ten longest layers. All hit the maximum $T_i =
1{,}835{,}008$ cycles and correspond to large FFN/attention
projection sublayers.}
\label{tab:longest_layers}
\begin{tabular}{rrl}
\toprule
\textbf{Layer idx} & \textbf{$T_i$ (cycles)} & \textbf{Selected config $c$} \\
\midrule
359 & 1{,}835{,}008 &  4 \\
376 & 1{,}835{,}008 &  4 \\
375 & 1{,}835{,}008 &  4 \\
368 & 1{,}835{,}008 &  4 \\
  0 & 1{,}835{,}008 &  4 \\
367 & 1{,}835{,}008 &  4 \\
718 & 1{,}835{,}008 &  0 \\
358 & 1{,}835{,}008 &  8 \\
357 & 1{,}835{,}008 &  8 \\
377 & 1{,}835{,}008 & 12 \\
\bottomrule
\end{tabular}
\end{table}

\paragraph{Placement ($p$, $p^W$): Global Buffer Residency}

\begin{itemize}[leftmargin=1.4em, itemsep=2pt]
    \item Activation edges pinned at $t=0$: $0\ /\ 703$
    \item Weight layers pinned at $t=0$: $0\ /\ 720$
    \item Total $p$-cells set across all $t$:
          $0\ /\ 506{,}160$ ($0.00\%$)
    \item Total $p^W$-cells set across all $t$:
          $1\ /\ 518{,}400$ ($0.00\%$)
\end{itemize}

\noindent The V1 LP makes almost no use of the global buffer for
long-lived tensor residency and nearly everything streams through.
With the L1/L2 buffer multipliers in this hardware configuration,
\textbf{tiling alone is sufficient to keep the chip fed}; persistent pins are not needed.

\paragraph{Rematerialization Events $r_{i,t}$}

\[
    \text{Layers with} \geq 1\ \text{remat event}: 0\ /\ 720.
    \qquad
    \text{Total remat events}: 0.
\]

\noindent Zero Rematerialization. V1 has the freedom to recompute
layers to free memory pressure, but on this cell determined it was
not worth the additional cycles. Activation streaming combined with
tile-config selection is sufficient.

\paragraph{Eviction / Restoration Timeline}

\begin{itemize}[leftmargin=1.4em, itemsep=2pt]
    \item Activation evictions / restorations: $0\ /\ 0$
    \item Weight evictions / restorations: $1\ /\ 1$
          (one weight churn across all 720 timesteps)
\end{itemize}

\paragraph{One Concrete Layer: Layer 360 (Deep Transformer Block)}

\[
    \text{Selected config}\ c = 0, \qquad T_i = 524{,}288\ \text{cycles}.
\]
\[
    \mathbf{y}_{\mathrm{bin}}\ \text{nonzero indices}: \{0,\ 16\}
\]
The LP relaxation assigned nonzero weight to two equally-good configs

\paragraph{What This Reveals About V1 on Llama-3 70B / H100}

\begin{enumerate}

    \item \textbf{Tile choice does most of the work, not placement or remat.} Of 1.6M decision variables, only ${\approx}720$
          (the $y_{i,c}$ selections) are doing real work; placement
          ($p$, $p^W$) and rematerialisation ($r$) variables are
          essentially all zero.

    \item \textbf{Large tiles win because H100 is large.} With
          1.05M MACs and 3.35 TB/s DRAM bandwidth, there is no memory
          slack pressure forcing small working sets. The LP exploits
          this by selecting the largest physically permissible tiles.

    \item \textbf{The $1.53\times$ lift over the naive baseline comes
          entirely from per-layer tile selection.} Ten different configs
          are active across the 720 layers, versus the naive
          one-size-fits-all heuristic. On smaller chips (e.g.\ TPU-v3)
          this same mechanism yields a $2.11\times$ lift, because the
          memory slack disappears and per-layer tile choice becomes
          correspondingly more impactful.

\end{enumerate}


% As a qualitative example of V1’s solution structure, we inspect the rounded V1 LP solution for DeepSeek-V3 on an RTX-5090-like hardware point. The model contains L=903 layers, including 174 MoE sublayers, and the V1 LP has 2.5M variables. Although the Pareto cache exposes 20 tile configurations per layer, the solution uses only five configs, {0,4,8,12,16}, suggesting that V1 compresses the available tiling space into a small workload-specific alphabet. The most revealing behavior occurs in the FFN sublayers: ffn\_up, ffn\_gate, and ffn\_down have identical compute time, yet V1 assigns them different tile configs in a rotating pattern. This indicates that the LP is not merely selecting the fastest per-layer tile; it is interleaving large- and small-working-set choices across adjacent sublayers to manage memory pressure. V1 also automatically identifies structural outliers: the only seven layers assigned the smallest tile c=16 correspond to unusual sublayers such as rotary-key projections, low-rank Q projections, shared/expert gates, and expert down-projections. Finally, weight churn is highly localized: only four attention-path weights in the first block are evicted/refetched, while the MoE sublayers stream without churn. This case illustrates how V1 exposes interpretable schedule structure that is unavailable in a sequential Timeloop→fusion pipeline.

\subsection{V1 LP Solution: DeepSeek-V3 on modeled RTX~5090}
This example shows the LP makes real per-sublayer scheduling decisions that are directly readable from layer names.


As a qualitative example of V1’s solution structure, we inspect the rounded V1 LP solution for DeepSeek-V3 on an RTX-5090-like hardware point. The model contains L=903 layers, including 174 MoE sublayers, and the V1 LP has 2.5M variables. Although the Pareto cache exposes 20 tile configurations per layer, the solution uses only five configs (Table~\ref{tab:ds_tile_configs}), {0,4,8,12,16}, suggesting that V1 compresses the available tiling space into a small workload-specific alphabet. The most revealing behavior occurs in the FFN sublayers: ffn\_up, ffn\_gate, and ffn\_down have identical compute time, yet V1 assigns them different tile configs in a rotating pattern (Table~\ref{tab:ds_ffn_stagger}). This indicates that the LP is not merely selecting the fastest per-layer tile; it is interleaving large- and small-working-set choices across adjacent sublayers to manage memory pressure. V1 also automatically identifies structural outliers: the only seven layers assigned the smallest tile c=16 correspond to unusual sublayers such as rotary-key projections, low-rank Q projections, shared/expert gates, and expert down-projections (Table~\ref{tab:ds_outliers}). Finally, weight churn is highly localized: only four attention-path weights in the first block are evicted/refetched (Table~\ref{tab:ds_evictions}), while the MoE sublayers stream without churn. This case illustrates how V1 exposes interpretable schedule structure that is unavailable in a sequential Timeloop→fusion pipeline.

\paragraph{Problem Dimensions}

\begin{itemize}[leftmargin=1.4em, itemsep=2pt]
    \item $L = 903$ layers (67-block DeepSeek-V3 with
          router/expert/shared sublayers per MoE block)
    \item $C = 20$ tile configs, $T = 903$ timesteps,
          $E = 886$ activation edges
    \item 2{,}504{,}019 LP variables,
          objective $= 517{,}356{,}013$ cycles
    \item Hardware (RTX~5090): 419K MACs, 1{,}792 GB/s DRAM,
          98 MB CGM
\end{itemize}

\paragraph{Tile-Config Selection: Only 5 of 20 Configs Ever Picked}

Configs are spaced exactly every 4 indices. V1 does not sweep the full continuum but locks onto a discrete alphabet of five shapes (Table~\ref{tab:ds_tile_configs}).

\begin{table}[ht!]
\centering
\caption{Tile-config distribution for DeepSeek-V3 on RTX~5090.
Only 5 of 20 candidate configs appear in the rounded solution.}
\label{tab:ds_tile_configs}
\begin{tabular}{rrl}
\toprule
\textbf{Config $c$} & \textbf{Layers} & \textbf{\% of model} \\
\midrule
0  & 363 & 40.2\% \quad (largest tile) \\
4  & 192 & 21.3\% \\
8  & 172 & 19.0\% \\
12 & 169 & 18.7\% \\
16 &   7 &  0.8\% \quad (smallest tile: tightest working set) \\
\midrule
others & 0 & unused \\
\bottomrule
\end{tabular}
\end{table}

\paragraph{Key Finding: V1 Staggers Tile Sizes Across the FFN Trio}

Every transformer block contains three FFN sublayers (up, gate, down) with identical compute time ($T_i = 2{,}580{,}482$ cycles each). V1 assigns \emph{different} tile configs to the three so their working sets do not peak simultaneously in the global buffer.
Table~\ref{tab:ds_ffn_stagger} shows the pattern across the ten hottest layers.

\begin{table}[ht!]
\centering
\caption{Top-10 longest layers. Despite identical $T_i$, V1 assigns
different tile configs to the FFN up/gate/down trio within each block, staggering working-set pressure across timesteps.}
\label{tab:ds_ffn_stagger}
\begin{tabular}{rllr}
\toprule
\textbf{Layer idx} & \textbf{Name} &
\textbf{$T_i$ (cycles)} & \textbf{Config $c$} \\
\midrule
 8 & \texttt{layer0\_ffn\_up}   & 2{,}580{,}482 &  0 \\
 9 & \texttt{layer0\_ffn\_gate} & 2{,}580{,}482 & 12 \\
10 & \texttt{layer0\_ffn\_down} & 2{,}580{,}482 &  8 \\
19 & \texttt{layer1\_ffn\_up}   & 2{,}580{,}482 &  4 \\
20 & \texttt{layer1\_ffn\_gate} & 2{,}580{,}482 &  0 \\
21 & \texttt{layer1\_ffn\_down} & 2{,}580{,}482 & 12 \\
30 & \texttt{layer2\_ffn\_up}   & 2{,}580{,}482 &  8 \\
31 & \texttt{layer2\_ffn\_gate} & 2{,}580{,}482 &  4 \\
32 & \texttt{layer2\_ffn\_down} & 2{,}580{,}482 &  0 \\
\bottomrule
\end{tabular}
\end{table}

This is the scheduling intelligence in V1: same compute cost,
different tiles, rotated so that any window of timesteps contains at most one large-working-set tile and two smaller ones. A naive mapper assigns one tile size to all FFN layers; V1 detects the working-set conflict and breaks the symmetry automatically.

\paragraph{Structural Outliers: The 7 Layers That Pick Config $c=16$}

The LP assigned the tightest working-set tile ($c = 16$) to exactly the sublayers with unusual shapes (rotary embeddings, shared-expert gates, and expert down-projections) without any hand-labeled outlier feature (Table~\ref{tab:ds_outliers}).

\begin{table}[h]
\centering
\caption{The seven layers assigned config $c=16$ (smallest tile /
tightest working set). All correspond to structurally atypical sublayers
in the DeepSeek-V3 MoE architecture.}
\label{tab:ds_outliers}
\begin{tabular}{rllr}
\toprule
\textbf{Layer idx} & \textbf{Name} &
\textbf{$T_i$ (cycles)} & \textbf{Reason} \\
\midrule
202 & \texttt{layer14\_K\_rope}      & 1{,}146{,}881 & Rotary-embedded K: small inner dim \\
238 & \texttt{layer16\_shared\_gate} &   286{,}720   & Shared-expert gate: narrow \\
322 & \texttt{layer22\_K\_rope}      & 1{,}146{,}881 & Same shape as layer 202 \\
394 & \texttt{layer27\_Q\_b}         &   737{,}281   & Q lower-rank decomposition \\
573 & \texttt{layer39\_Q\_a}         &   215{,}040   & Q upper-rank decomposition \\
706 & \texttt{layer47\_expert\_gate} &   286{,}720   & Per-expert gate \\
812 & \texttt{layer54\_expert\_down} &   286{,}720   & Per-expert down-projection \\
\bottomrule
\end{tabular}
\end{table}

\paragraph{Weight Evictions: Localized to Transformer Block~1}

V1 evicts weights for only four layers, all in the first transformer block's attention path (Table~\ref{tab:ds_evictions}).

\begin{table}[h]
\centering
\caption{Weight eviction/restoration events. Only transformer block~1
attention weights are evicted; all 174 MoE blocks have zero churn.}
\label{tab:ds_evictions}
\begin{tabular}{rlrrrr}
\toprule
\textbf{Layer} & \textbf{Name} &
\textbf{Evict} & \textbf{Restore} &
\textbf{$T_i$ (cycles)} & \textbf{Config $c$} \\
\midrule
12 & \texttt{layer1\_Q\_b}   & 1 & 0 & 737{,}281   & 0  \\
13 & \texttt{layer1\_KV\_a}  & 1 & 1 &  71{,}680   & 12 \\
14 & \texttt{layer1\_KV\_b}  & 1 & 1 & 327{,}680   &  8 \\
15 & \texttt{layer1\_K\_rope}& 1 & 1 & 1{,}146{,}881 & 4 \\
\bottomrule
\end{tabular}
\end{table}

\noindent The LP determined that block~0/1 attention weights are cheaper to refetch than to evict the FFN trio they compete with in the global buffer.

\paragraph{Relative Layer Costs}
Table~\ref{tab:ds_layer_costs} summarizes the execution time of key sublayer types.
\begin{table}[h]
\centering
\caption{Relative execution time of key sublayer types. MoE routers
are ${\approx}72\times$ cheaper than FFN sublayers yet still receive varied tile assignments because 67 of them serialize in the dependency chain.}
\label{tab:ds_layer_costs}
\begin{tabular}{lrr}
\toprule
\textbf{Sublayer type} & \textbf{$T_i$ (cycles)} &
\textbf{Ratio vs.\ FFN} \\
\midrule
FFN up / gate / down & 2{,}580{,}482 & $1.00\times$ (long pole) \\
Q/K/V projections    & 700K–1{,}200K & ${\approx}27$–$47\%$ \\
MoE routers          & 35{,}840      & $1.4\%$ ($72\times$ smaller) \\
\bottomrule
\end{tabular}
\end{table}

\paragraph{What This Demonstrates About V1}

\begin{enumerate}

    \item \textbf{V1 performs scheduling, not just per-layer tile
          picking.} The FFN trio tile rotation is a genuine scheduling decision: no naive per-layer mapper can produce this.

    \item \textbf{V1 discovers structural outliers automatically.}
          The seven $c=16$ picks land exactly on K\_rope, Q\_a, Q\_b,
          expert\_gate, and expert\_down sublayers without any
          hand-labeled outlier feature in the input.

    \item \textbf{V1 commits to a 5-config alphabet, not a 20-config
          sweep.} This is actionable for hardware designers: only 5 of
          20 candidate tile shapes ever matter for this workload.

    \item \textbf{Memory churn is highly localized.} Only 4 evictions
          occur, all in transformer block~1, all in the attention path.
          The 174 MoE blocks have zero churn: V1 streams them entirely.

\end{enumerate}

\begin{figure}[ht!]
    \centering
    \includegraphics[width=0.7\linewidth]{figs4/fft2.png}
    \caption{Visualization or rotating tiles selected by V1 LP. There exists 903 FFN sublayers in this workload, tile config C=20. The LP only picks 5 of them (since the rest are Pareto-dominated).}
    \label{fig:tiling_rotate}% RENAMED: was a duplicate of solver_app.tex's fig:pipeline, which silently broke \ref{fig:pipeline} in prelim.tex
\end{figure}
\label{app:tiling_demo}% RENAMED: was a duplicate of solver_app.tex's app:vizier, which silently broke \ref{app:vizier} in prelim.tex